\chapter{Tests and Results}
\label{ch:tests-and-results}

No scientific runtime result is claimed from the failed Leonardo launch because
neither the diagnostic executable nor BRINE began. This chapter therefore
retains the planned validation structure without inventing measurements.

\section{Evidence Status}

Table~\ref{tab:evidence-status} defines the distinction used when results are
added. It prevents source inspection, prepared inputs, and successful runtime
checks from being reported as equivalent forms of evidence.

\begin{table}[htbp]
  \centering
  \caption{Evidence classifications used for the project report.}
  \label{tab:evidence-status}
  \begin{spacing}{1.0}
  \small
  \begin{tabularx}{\textwidth}{@{}>{\bfseries}lX@{}}
    \toprule
    Classification & Meaning in this report \\
    \midrule
    Source-confirmed & Established from the active implementation, immutable
      input, or a directly identified reference artifact. \\
    Runtime-confirmed & Demonstrated by a completed run and its stated
      acceptance checks. \\
    Not tested & No relevant successful runtime observation is available. \\
    Blocked & A failed prerequisite or absent, unapproved implementation
      prevents the test. \\
    \bottomrule
  \end{tabularx}
  \end{spacing}
\end{table}

\section{Planned Validation Sequence}

The completed report will present deterministic boundary checks, canonical
fresh-start behaviour, checkpoint and restart continuity, time-step evidence,
one-node performance, and two-node launch and scaling results only when the
corresponding runtime evidence exists. Comparison plots will follow the
project's established scientific plotting conventions.
